\documentclass[11pt]{article}
\usepackage[margin=1in]{geometry}
\usepackage{booktabs}
\usepackage{graphicx}
\usepackage{amsmath}
\usepackage{amssymb}
\usepackage{float}
\usepackage{array}
\title{Tip or Skip: Uncertainty-Aware NYC Taxi Tipping Prediction}
\author{Wali Ahmed \and Geoffrey Kim \and Jiachen Tu}
\date{May 2026}
\begin{document}
\maketitle

\begin{abstract}
This project investigates recorded electronic tipping behavior in New York City taxi trips using official TLC Yellow and Green Taxi records. The central hypothesis is that tipping should be modeled as a calibrated two-stage distributional problem rather than a single point-regression problem: first predict whether a trip receives any recorded electronic tip, then model the conditional distribution of positive tips. We build a frozen 2024--2025 dataset with 1,008,000 cleaned credit-card trips, compare logistic/ridge and boosted-tree hurdle baselines against a Tabular Transformer Mixture Density Network (MDN), and deploy the results in an interactive Hugging Face demo with a grounded tipping-facts assistant, what-if sensitivity sweeps, model comparison panels, maps, and a risk-aware shift planner. The strongest point-prediction model is the tree hurdle baseline, with expected-tip MAE of \$1.45, while the Transformer-MDN provides uncertainty estimates and lower-tail quantiles that support decision-focused analysis.
\end{abstract}

\section{Introduction}
Taxi tipping is a small individual decision repeated millions of times across a dense urban transportation system. It is also a useful applied machine learning problem because the observed target is skewed, partially censored, and shaped by trip context. A naive formulation would ask for one number: the predicted tip amount. That formulation hides two important facts. First, a trip may or may not receive a recorded electronic tip at all. Second, when a tip occurs, the amount can vary widely across airport trips, short neighborhood trips, boroughs, time of day, and fare levels.

Our project asks whether a hurdle-style, uncertainty-aware formulation produces a more useful model than a single point predictor. The practical application is a driver-facing or analyst-facing decision-support tool: compare zones, inspect risk, ask natural-language questions about the data, and run what-if experiments. The machine learning question is broader: when does a deep distributional model add value on structured tabular data where strong classical baselines are already available?

This framing follows the project guidelines by focusing on a testable hypothesis and model behavior rather than only leaderboard performance. We include strong baselines, a deep generative component, ablation through model-family comparisons, subgroup evaluation, and a discussion of negative results. The final result is not simply ``train a taxi model.'' It is a small research-style study of calibration, uncertainty, and decision-support outputs for a real tabular dataset.

The project is also designed around inspectability. A useful applied model should let a user ask why one model is preferred, where the model is reliable, and how predictions move when inputs change. For that reason, the final artifact includes both a technical report and an interactive demo. The demo is not separate from the experiment; it is a way to expose model behavior, subgroup performance, and uncertainty to a reader who wants to interrogate the results directly.

\section{Related Work and Motivation}
The modeling design draws on three ideas. The first is hurdle modeling: when outcomes have a mass at zero and a continuous positive tail, separating occurrence from amount can be more interpretable and often easier to optimize than fitting one regressor. Tipping naturally has this structure because the event of leaving any recorded electronic tip is distinct from the conditional amount.

The second idea is strong tabular baselines. Gradient-boosted tree models remain difficult to beat on many structured datasets because they capture nonlinear threshold effects and feature interactions without requiring large-scale representation learning. For that reason, the project treats boosted trees as a serious baseline rather than as a strawman.

The third idea is distributional prediction. Mixture Density Networks model a conditional density instead of only a mean. In this project the MDN head estimates a mixture distribution over positive log tips. That makes the deep model useful even when its point-prediction error is not the best result: it gives quantiles, interval coverage, and lower-tail risk estimates that can be used directly in the shift planner.

A fourth motivation is calibration. In a decision-support setting, a probability should be meaningful, not just rank examples correctly. A driver comparing two options needs to know whether a predicted probability and expected value are stable enough to trust. We therefore report both discrimination metrics such as ROC-AUC and calibration-sensitive metrics such as Brier score and expected calibration error.

\section{Dataset and Preprocessing}
The dataset is built from official NYC TLC Yellow and Green taxi monthly parquet files for 2024 and 2025. We retain credit-card trips only because TLC \texttt{tip\_amount} records electronic tips and does not include cash tips. This choice avoids mixing observed zeros with unobserved cash tips, but it also means the target is recorded electronic tipping behavior rather than all real-world tipping.

Cleaning removes rows with nonpositive fare, distance, or duration, invalid pickup-month alignment, and inconsistent basic trip fields. We derive pickup year, month, hour, weekday, weekend indicator, daypart, trip duration in minutes, passenger-count buckets, rate-code categories, and normalized store-and-forward flags. Pickup and dropoff location identifiers are joined to TLC taxi-zone names, boroughs, and service zones. The final model feature set contains temporal, fare, surcharge, location, taxi-type, vendor, and trip-structure variables.

The frozen dataset contains 1,008,000 rows and 46 columns. Its overall recorded electronic tip rate is 92.9%. The split is chronological: January--September 2024 for training (378,000 rows), October--December 2024 for validation (126,000 rows), and all of 2025 for testing (504,000 rows). This setup makes the held-out test split a future-period evaluation instead of a random split that could overstate performance through temporal leakage.

We also remove leakage-prone variables from the model inputs. In particular, \texttt{total\_amount} is excluded because it can encode the tip after the fact, and the task is restricted to credit-card trips rather than letting payment type directly reveal target observability. The retained variables describe trip context: time, distance, duration, fare components, taxi type, vendor, passenger bucket, rate code, and pickup/dropoff geography. This matches the interactive prediction setting, where the user asks what the model expects from a hypothetical trip before the tip is known.

\section{Machine Learning Formulation}
For a trip $i$ with features $x_i$, Stage 1 defines
\[
y_i = \mathbb{1}[\texttt{tip\_amount}_i > 0]
\]
and estimates $P(y_i=1 \mid x_i)$. Stage 2 models the transformed positive amount
\[
z_i = \log(1+\texttt{tip\_amount}_i)
\]
for trips where $y_i=1$. The expected electronic tip is then approximated by combining the probability of a tip with the conditional positive-tip prediction.

For the MDN model, the conditional positive-tip density is
\[
p_\theta(z \mid y=1,x)=\sum_{k=1}^K \pi_k(x)\mathcal{N}(z;\mu_k(x),\sigma_k^2(x)).
\]
This formulation supports multiple summaries of the same trip: expected tip, median tip, lower-tail quantiles, and interval coverage. We evaluate classification with ROC-AUC, average precision, log loss, Brier score, expected calibration error, precision, recall, and F1. We evaluate conditional amount prediction with log-tip MAE/RMSE and evaluate full decision output with expected-tip MAE.

\section{Models and Ablations}
We compare three model families. The logistic-ridge baseline combines logistic regression for Stage 1 with ridge regression for Stage 2. It tests how much signal is recoverable from a simple linear model with regularization. The tree hurdle baseline uses histogram gradient-boosted trees for both stages. It is the strongest classical tabular baseline and captures nonlinear interactions among fare, distance, time, and zone variables.

The deep model embeds categorical features and combines them with normalized numeric features in a Tabular Transformer. It has a binary classification head and an MDN head for the positive-tip distribution. This is the project requirement's deep/generative component. The comparison across linear, tree, and Transformer-MDN models functions as the main ablation: it asks whether the extra representational and distributional machinery improves point prediction, uncertainty, or decision-support outputs.

The ablation is intentionally framed around model behavior. The linear baseline tests whether the engineered features are already sufficient under a simple decision boundary. The tree model tests whether nonlinear tabular interactions explain most of the signal. The Transformer-MDN tests whether a deep representation and a generative output distribution add value beyond those tabular baselines. This lets us interpret the final outcome even when the deep model does not win every metric.

\section{Experimental Results}
\begin{table}[H]
\centering
\caption{Final model comparison on the held-out 2025 test split.}
\resizebox{\linewidth}{!}{%
\begin{tabular}{lrrrrrrrr}
\toprule
Model & ROC-AUC & Log loss & Brier & ECE & F1 & Log-tip RMSE & Expected-tip MAE & 80\% cov. \\
\midrule
logistic\_ridge & 0.737 & 0.574 & 0.191 & 0.354 & 0.859 & 0.339 & \$2.51 & -- \\
tree\_hurdle & 0.771 & 0.205 & 0.052 & 0.004 & 0.969 & 0.319 & \$1.45 & -- \\
transformer\_mdn & 0.732 & 0.581 & 0.196 & 0.366 & 0.867 & 0.326 & \$2.34 & 0.743 \\
\bottomrule
\end{tabular}}
\end{table}

\begin{figure}[H]
\centering
\includegraphics[width=\linewidth]{figures/model_comparison.png}
\caption{Comparison across classification, calibration, conditional amount, and expected-tip metrics.}
\end{figure}

The tree hurdle model is the best point predictor, achieving the lowest expected-tip MAE. This is a meaningful negative result for the deep model: the Tabular Transformer-MDN does not dominate a strong boosted-tree baseline on structured taxi data. However, the deep model changes the form of the output. It provides an estimated conditional distribution, which allows the project to rank zones using lower-tail risk and to report interval coverage. This is more useful for the interactive shift-planning setting than a single regression output.

The results therefore separate two notions of success. For point prediction, the boosted-tree hurdle model is the most practical choice. For uncertainty-aware analysis, the Transformer-MDN contributes information the tree model does not directly provide. This distinction is the main insight of the project: the best model for a leaderboard metric is not necessarily the only model worth deploying in an exploratory ML interface.

\begin{figure}[H]
\centering
\includegraphics[width=\linewidth]{figures/monthly_tip_rate.png}
\caption{Recorded electronic tip rate by month for Yellow and Green taxi trips in the frozen dataset.}
\end{figure}


\section{Ablation, Calibration, Temporal, Graph, and Language Experiments}
After the main model comparison, the experiment path turns to model behavior. The next set of experiments asks what information the model uses, whether predicted probabilities are calibrated, whether taxi-zone flow structure helps, whether recent hourly history predicts the next hour, and whether the driver-facing language layer stays grounded in numeric model outputs. These experiments are part of the same evaluation pipeline as the baselines and Transformer-MDN, because they test reliability and usability rather than only point error.

\begin{table}[H]
\centering
\caption{Feature ablation metrics. Lower expected-tip MAE and ECE are better; higher ROC-AUC is better.}
\begin{tabular}{lrrrr}
\toprule
Ablation & Features & ROC-AUC & ECE & Expected-tip MAE \\
\midrule
no\_time & 24 & 0.765 & 0.004 & \$1.46 \\
full & 30 & 0.769 & 0.004 & \$1.46 \\
no\_distance & 28 & 0.767 & 0.004 & \$1.46 \\
no\_zones & 22 & 0.750 & 0.004 & \$1.47 \\
no\_fare & 21 & 0.769 & 0.004 & \$1.47 \\
\bottomrule
\end{tabular}
\end{table}

\begin{table}[H]
\centering
\caption{Calibration bins for the selected hurdle model.}
\begin{tabular}{lrrrr}
\toprule
Probability bin & Rows & Predicted & Actual & Gap \\
\midrule
0.0--0.1 & 9,194 & 0.000 & 0.000 & -0.000 \\
0.1--0.2 & 3 & 0.153 & 0.000 & -0.153 \\
0.2--0.3 & 20 & 0.272 & 0.200 & -0.072 \\
0.3--0.4 & 113 & 0.358 & 0.301 & -0.057 \\
0.4--0.5 & 490 & 0.463 & 0.520 & 0.057 \\
0.5--0.6 & 1,450 & 0.554 & 0.548 & -0.005 \\
0.6--0.7 & 2,568 & 0.656 & 0.662 & 0.006 \\
0.7--0.8 & 9,301 & 0.760 & 0.763 & 0.003 \\
0.8--0.9 & 36,250 & 0.862 & 0.858 & -0.004 \\
0.9--1.0 & 444,611 & 0.958 & 0.954 & -0.004 \\
\bottomrule
\end{tabular}
\end{table}


\begin{figure}[H]
\centering
\includegraphics[width=0.82\linewidth]{figures/ablation_expected_tip_mae.png}
\caption{Feature ablation results using expected-tip MAE.}
\end{figure}


\begin{figure}[H]
\centering
\includegraphics[width=0.62\linewidth]{figures/calibration_curve.png}
\caption{Calibration curve comparing predicted tip probability with observed tip rate.}
\end{figure}


The ablation answers a basic question: what is the model leaning on? Removing zone features, time features, fare fields, or distance fields changes the error profile and makes the strongest feature groups visible. The calibration curve answers a different question: when the model says that a tip is likely, does that probability behave like a probability? These checks are important because the demo uses predicted probabilities to rank rides.

The graph-flow experiment represents TLC zones as nodes and pickup-dropoff traffic as edges. Node features include pickup volume, dropoff volume, in-degree, out-degree, average tip, and tip rate from the training period. A small graph convolution model then predicts future zone-level tip behavior. The graph model used 259 TLC zone nodes. Its test zone-tip MAE was \$1.21; the temporal naive baseline was \$0.41. This was a useful negative result because it showed that simple year-to-year zone stability was hard to beat.


\begin{figure}[H]
\centering
\includegraphics[width=0.82\linewidth]{figures/zone_flow_graph_summary.png}
\caption{Busiest taxi-zone nodes in the held-out pickup-dropoff flow graph.}
\end{figure}


\begin{figure}[H]
\centering
\includegraphics[width=0.68\linewidth]{figures/graph_gcn_zone_prediction.png}
\caption{Graph model zone-tip predictions against 2025 observed zone averages.}
\end{figure}


The temporal sequence experiment aggregates rides into hourly time series by taxi type and pickup borough. A small LSTM reads the previous six hourly summaries and predicts the next hour's tip rate and average tip. The LSTM next-hour experiment used 43,065 test windows. Its average-tip MAE was \$1.80, compared with \$2.40 for a last-hour naive baseline. The tip-rate MAE was 0.148. This connects the project to the RNN/LSTM part of the course and gives a first version of shift-planning over time rather than only over zones.


\begin{figure}[H]
\centering
\includegraphics[width=0.75\linewidth]{figures/sequence_lstm_training.png}
\caption{Sequence LSTM training and validation loss.}
\end{figure}


\begin{figure}[H]
\centering
\includegraphics[width=0.62\linewidth]{figures/sequence_shift_signal.png}
\caption{Current-hour versus next-hour tip-rate signal.}
\end{figure}


The driver-copilot layer was also evaluated directly. We created scripted ride-choice questions with known expected zones and expected dollar values from the final zone-risk table, then checked whether the answer included the recommended zone, the expected-tip number, and the cash-tip limitation. The scripted driver-copilot evaluation reached an overall grounding score of 100.0%. The same zone-risk table was converted into instruction-tuning examples for a compact driver LLM. The compact driver LLM run fine-tuned Qwen/Qwen2.5-0.5B-Instruct on 161 generated ride-planning examples, with eval perplexity 1.35 and zone mention rate 100.0%. The live demo still keeps deterministic retrieval as the default because it is more reliable for grading, but the fine-tuning run shows how the project can support a specialized driver-facing language model.


\begin{figure}[H]
\centering
\includegraphics[width=0.62\linewidth]{figures/copilot_eval_summary.png}
\caption{Driver Copilot grounding check pass rates.}
\end{figure}



\section{Subgroup Behavior and Spatial Risk}
Because the dataset is geographically uneven, aggregate metrics can hide important differences. We therefore compute subgroup metrics by taxi type and pickup borough. Large Manhattan slices dominate the data, while smaller borough slices tend to be noisier and less calibrated. This matters for deployment because a driver-facing recommendation should show uncertainty and observed trip counts, not only a sorted list of high-value zones.

\begin{table}[H]
\centering
\caption{Largest Transformer-MDN subgroup evaluation slices.}
\begin{tabular}{llrrr}
\toprule
Taxi & Pickup borough & Rows & ROC-AUC & F1 \\
\midrule
yellow & Manhattan & 319,329 & 0.632 & 0.960 \\
green & Manhattan & 95,085 & 0.693 & 0.741 \\
yellow & Queens & 33,615 & 0.786 & 0.769 \\
green & Queens & 27,821 & 0.632 & 0.081 \\
green & Brooklyn & 20,035 & 0.579 & 0.291 \\
yellow & Brooklyn & 4,777 & 0.983 & 0.551 \\
\bottomrule
\end{tabular}
\end{table}

\begin{table}[H]
\centering
\caption{Top final-model pickup zones by expected tip among zones with at least 50 observed test rows.}
\resizebox{\linewidth}{!}{%
\begin{tabular}{lllrrr}
\toprule
Taxi & Borough & Zone & Rows & Expected tip & Q10 tip \\
\midrule
yellow & Queens & JFK Airport & 16,481 & \$6.15 & \$6.38 \\
yellow & Queens & LaGuardia Airport & 12,815 & \$5.84 & \$5.66 \\
yellow & Queens & East Elmhurst & 1,244 & \$4.29 & \$5.45 \\
yellow & Queens & Flushing Meadows-Corona Park & 182 & \$3.99 & \$5.00 \\
green & Brooklyn & Red Hook & 606 & \$3.49 & \$4.37 \\
green & Brooklyn & Flatbush/Ditmas Park & 72 & \$3.47 & \$5.12 \\
\bottomrule
\end{tabular}}
\end{table}

\begin{figure}[H]
\centering
\includegraphics[width=0.9\linewidth]{figures/top_zone_expected_tip.png}
\caption{Top Manhattan pickup zones by predicted expected tip according to the final uncertainty-aware model.}
\end{figure}

\section{Interactive Hugging Face Demo}
The Hugging Face Space is part of the final ML artifact, not only a visualization wrapper. The demo includes a grounded tipping-facts assistant that answers questions using the frozen dataset summary, final metrics, subgroup table, and zone-risk table. If a Hugging Face Inference API token and model are configured, the app can route these grounded facts through a hosted language model; otherwise it uses deterministic artifact retrieval. This keeps the assistant reliable for grading and prevents unsupported claims.

The demo also includes a prediction form for hypothetical credit-card trips, a what-if sensitivity sweep over pickup hour, fare, distance, or duration, final model-comparison controls, monthly dataset profiles, subgroup metrics, a Manhattan map, and a risk-aware shift planner. The shift planner can optimize for expected tip, downside $Q_{0.10}$ tip, or predicted tip probability. These controls expose the machine learning behavior directly: users can see how changing inputs changes predictions and how different objectives change zone rankings.

The intended inspection workflow is sequential. A reader can first ask the assistant what dataset and target were used, then compare model metrics in the model lab, then construct a hypothetical trip and run a sensitivity sweep, and finally inspect whether zone recommendations change under risk-neutral versus risk-averse objectives. This makes the demo a compact model-audit tool rather than a static dashboard.

\section{Driver-Facing LLM Copilot}
The final interface includes a driver-facing LLM layer called Driver Copilot. Its purpose is to let a driver ask natural questions about ride choice and shift planning, such as: ``I am at Midtown Center and got two ride options, JFK Airport or LaGuardia Airport. Which one should I choose?'' The copilot parses the prompt, identifies known TLC zones and aliases, retrieves final model outputs for those zones, and answers with a recommendation grounded in expected tip, downside $Q_{0.10}$ tip, predicted tip probability, and observed held-out trip count.

The live demo uses retrieval-grounded answers as the default because the core evidence is structured and numeric. The deterministic layer maps driver language to zone-level and trip-level model outputs. We also generated instruction-tuning examples from the zone-risk table and trained a compact driver LLM on those examples. The tuned LLM is evaluated separately for grounding behavior, while the public Space keeps the deterministic layer available even when no Hugging Face inference token is configured.

This component turns the machine learning results into a usable driver workflow. For a single area, the copilot classifies the area as strong, solid, or lower relative to comparable zones. For two candidate areas, it recommends the option with higher expected electronic tip and reports the gap. The structured comparison form additionally uses the deployed two-stage trip predictor to compare two rides with specified pickup area, dropoff areas, hour, weekday, month, distance, fare, and duration. The result is an LLM-style planning layer whose outputs are auditable rather than free-form.

\section{Limitations and Ethics}
The most important limitation is target observability. Cash tips are not recorded in TLC \texttt{tip\_amount}, so the model should be described as predicting recorded electronic tips. A zero in the data does not necessarily mean a rider left no tip; it means no electronic tip was recorded. This affects interpretation, especially across neighborhoods or trip types where cash behavior may differ.

There are also deployment risks. Zone rankings could amplify existing demand patterns if treated as instructions rather than analysis. Small-sample zones can produce unstable estimates, so the demo exposes observed trip counts and lower-tail quantiles. The model should be used as decision support, not as a guarantee of income or as a replacement for local driver knowledge.

The model is not causal. It can show that certain trip contexts are associated with higher recorded electronic tips, but it cannot prove that choosing a particular route will cause a higher tip for a particular driver. Passenger mix, traffic, airport rules, event timing, and unobserved rider behavior all affect outcomes. This is why the final interface emphasizes comparison, uncertainty, and evidence rather than prescriptive instructions.

\section{Conclusion}
Tip or Skip reframes NYC taxi tipping prediction as an uncertainty-aware applied ML problem. The project produces a reproducible dataset, clear ML formulation, strong baselines, a deep generative component, ablations through model comparisons, subgroup and spatial analyses, and an interactive demo. The main empirical finding is pragmatic: boosted trees are the best point predictor, while the Transformer-MDN contributes distributional information that makes the system more useful for risk-aware exploration. This gives a more honest and more interesting result than simply reporting an accuracy number.

\section{Reproducibility}
The frozen dataset package is stored under \texttt{Prototype/final\_dataset}. The main training and reporting commands are:
\begin{verbatim}
python scripts/train_baselines.py
python scripts/train_transformer_mdn.py --epochs 16 --batch-size 8192
python scripts/evaluate_models.py
python scripts/generate_latex_report.py
\end{verbatim}
The final project materials include this PDF, the offline \texttt{index.html} technical blog, the Hugging Face Space source, plots, metrics, model artifacts, and the frozen dataset package.

\begin{thebibliography}{9}
\bibitem{tlc} New York City Taxi and Limousine Commission. TLC Trip Record Data and Data Dictionaries.
\bibitem{mdn} Bishop, C. M. Mixture Density Networks. Aston University technical report, 1994.
\bibitem{tabtransformer} Huang, X. et al. TabTransformer: Tabular Data Modeling Using Contextual Embeddings. 2020.
\bibitem{calibration} Guo, C. et al. On Calibration of Modern Neural Networks. ICML, 2017.
\bibitem{gcn} Kipf, T. N. and Welling, M. Semi-Supervised Classification with Graph Convolutional Networks. ICLR, 2017.
\bibitem{lstm} Hochreiter, S. and Schmidhuber, J. Long Short-Term Memory. Neural Computation, 1997.
\bibitem{lora} Hu, E. J. et al. LoRA: Low-Rank Adaptation of Large Language Models. 2021.
\end{thebibliography}

\end{document}
